\documentclass[10pt,twocolumn,twoside,slovak,a4paper]{article}

\usepackage[T1]{fontenc}
\usepackage[utf8]{inputenc}
\usepackage[slovak]{babel}
\usepackage{graphicx}
\usepackage{wrapfig}
\usepackage{amsmath}
\usepackage{lipsum}
\usepackage{mwe} % poskytuje example-image-a/b/c (dočasné obrázky)
\graphicspath{{./}{images/}} % ak máš obrázky v "images/", nechaj to tak

\title{Predbežná verzia článku – Cvičenie 4}
\author{Michal Pípa, Michal Urban, Tomáš Topalov, Tomáš Šuba}
\date{Ak. rok 2025/26}

\begin{document}
\maketitle

\begin{abstract}
Dokument k cvičeniu č. 4 (dvojstĺpcový layout, dva diagramy vedľa seba, obrázok cez oba stĺpce, logo FIIT obtekané textom, matica a vzorec).
\end{abstract}

% Logo FIIT obtekané textom (dočasne použijeme example-image-a)
\begin{wrapfigure}{r}{0.28\textwidth}
  \vspace{-8pt}
  \centering
  \includegraphics[width=0.28\textwidth]{example-image-a} % nahraď za logo_fiit.png
  \vspace{-8pt}
\end{wrapfigure}

\section{Úvod}
\lipsum[1]

% Dva diagramy vedľa seba (dočasne example-image-b/c)
\begin{figure}[h]
\centering
\includegraphics[width=0.48\linewidth]{example-image-b}%
\hfill
\includegraphics[width=0.48\linewidth]{example-image-c}
\caption{Dva UML diagramy vedľa seba. (nahraď za diagram1/diagram2)}
\end{figure}

\lipsum[2]

% Obrázok cez oba stĺpce
\begin{figure*}[t]
\centering
\includegraphics[width=\textwidth]{example-image} % nahraď za diagram_fullwidth.png
\caption{Obrázok cez oba stĺpce.}
\end{figure*}

\section{Matica a vzorec}
\[
A=
\begin{bmatrix}
1 & 2 & 3 & 4 \\
5 & 6 & 7 & 8 \\
9 & 10 & 11 & 12 \\
13 & 14 & 15 & 16 \\
17 & 18 & 19 & 20
\end{bmatrix}
\]

\[
S_n = \sum_{k=1}^{n} \frac{1}{k} + \frac{1}{n+1} + \frac{1}{n+2} + \cdots
\]

\section{Porovnanie s MS Word}
V \LaTeX{} sa matice a vzorce píšu značkami (vyššia presnosť a znovupoužiteľnosť); vo Worde je to vizuálne rýchlejšie, no menej flexibilné pri zložitejších prípadoch.

\section{Záver}
Dokument spĺňa požiadavky Cvičenia~4.

\end{document}
